%% FullCourse_10Lecs -- Lecture 9, Part 9.2b
%% Assembled by tools/lec09_assemble.py from reorg_manifest.yaml: Phase A of
%% Lec09_Reorg_Plan_2026-09-12.md as amended by Lec09_Reorg_Plan_Review_2026-09-12.md.
%% Each "% ---- <ID> · <TAG> · TIER" marker names a frame's source deck and line;
%% keep the markers until Phase B is done (lec09_assemble.py check reads them).
%% Owns: terminal, install, Git, global vs project files, the harness file mapping (verified, review A2.5)
%% Owns: initialization, three memories, the resolution rule, the twelve commands, modes, keys, permissions
%% Owns: first session, tokens, the context window, compaction mechanics
%% Defers: rung surfaces (/agents, --agent, --agents, -p, --bg) -> T3b; dated install detail -> TA; hooks and settings.json -> T5b

%% Shared preamble for the FullCourse_10Lecs topic decks.
%%
%% Each topic .tex \input's this file, then defines its own \title, \date
%% override (if any), \graphicspath, and \begin{document}.
%%
%% Reconstructed verbatim from the inlined copy preserved in
%% Lec10_Case_Studies/Lec10_T8_Case6_DeepRamsey.tex (lines 23-190), which is
%% explicitly annotated there as "from ../shared_preamble.tex, copied verbatim".
%% Base preamble matches MiniCourse_8hr/Slides/shared_preamble.tex; this version
%% additionally provides \sectiondivider, the conceptbox/cautionbox/examplebox/
%% takeawaybox tcolorboxes, and \compacttables used across the split decks.

\documentclass[10pt,english,aspectratio=169]{beamer}

\usepackage{ae,aecompl}
\usepackage[T1]{fontenc}
\usepackage{textcomp}
\usepackage[utf8]{inputenc}
\usepackage{babel}
\usepackage{datetime2}

\setcounter{secnumdepth}{3}
\setcounter{tocdepth}{3}

\usepackage{amsmath, amssymb, amsthm, graphicx}
\usepackage{booktabs, multirow}
\usepackage{tikz}
\usepackage{pgfplots}
\pgfplotsset{compat=1.16}
\usepackage[most]{tcolorbox}
\tcbuselibrary{skins, breakable, theorems}
\usepackage{listings}
\usepackage{xcolor}

\lstset{
  basicstyle=\ttfamily\small,
  keywordstyle=\color{blue!80!black}\bfseries,
  commentstyle=\color{green!50!black}\itshape,
  stringstyle=\color{orange!80!black},
  backgroundcolor=\color{gray!8},
  frame=single,
  rulecolor=\color{gray!40},
  breaklines=true,
  showstringspaces=false,
  numbers=none,
}

\lstdefinestyle{pythonstyle}{
    language=Python,
    basicstyle=\ttfamily\footnotesize,
    keywordstyle=\color{blue!80!black}\bfseries,
    commentstyle=\color{green!50!black}\itshape,
    stringstyle=\color{orange!80!black},
    frame=single,
    breaklines=true,
    showstringspaces=false,
    numbers=left,
    numbersep=5pt,
    numberstyle=\tiny\color{gray},
    backgroundcolor=\color{gray!8},
    rulecolor=\color{gray!40}
}

\ifx\hypersetup\undefined
  \AtBeginDocument{%
    \hypersetup{bookmarks=false,breaklinks=true,pdfborder={0 0 0},colorlinks=true,
      linkcolor=DarkText,urlcolor=RedTitle}%
  }
\else
  \hypersetup{bookmarks=false,breaklinks=true,pdfborder={0 0 0},colorlinks=true,
    linkcolor=DarkText,urlcolor=RedTitle}%
\fi

\makeatletter
\newcommand\makebeamertitle{%
  \begin{frame}[plain]
    \vspace*{1.0cm}
    \centering
    {\usebeamerfont{title}\usebeamercolor[fg]{title}\inserttitle\par}
    \vspace{1.0cm}
    {\usebeamerfont{author}\usebeamercolor[fg]{author}\insertauthor\par}
    \vspace{0.2cm}
    {\usebeamerfont{date}\usebeamercolor[fg]{date}\insertdate\par}
  \end{frame}}%
\AtBeginDocument{%
  \let\origtableofcontents=\tableofcontents
  \def\tableofcontents{\@ifnextchar[{\origtableofcontents}{\gobbletableofcontents}}
  \def\gobbletableofcontents#1{\origtableofcontents}%
}

\usetheme{metropolis}
\usefonttheme{professionalfonts}

\definecolor{LightGray}{RGB}{230,230,230}
\definecolor{RedTitle}{RGB}{0,0,200}
\definecolor{VeryLightGray}{RGB}{245,245,245}
\definecolor{DarkText}{RGB}{0,0,0}

\setbeamercolor{background canvas}{bg=white}
\setbeamercolor{title}{fg=RedTitle}
\setbeamercolor{author}{fg=DarkText}
\setbeamercolor{date}{fg=DarkText}
\setbeamercolor{frametitle}{bg=VeryLightGray,fg=DarkText}
\setbeamercolor{normal text}{fg=DarkText}
\setbeamerfont{title}{size=\large}

\setbeamersize{text margin left=5mm, text margin right=5mm}
\addtobeamertemplate{frametitle}{}{\vspace{-0.5em}}
\newcommand{\reducevspace}{\vspace{-0.5cm}}

% Shared section divider and box styles for split topic decks.
\newcommand{\sectiondivider}[2]{%
\begin{frame}[plain,noframenumbering]
\begin{center}
\vfill
{\Large\textbf{#1}\par}
\vspace{0.35cm}
{\normalsize #2\par}
\vfill
\end{center}
\end{frame}
}

\newtcolorbox{conceptbox}[1][]{
  colback=blue!3,
  colframe=RedTitle!55,
  boxrule=0.35mm,
  arc=1mm,
  left=2mm,
  right=2mm,
  top=1mm,
  bottom=1mm,
  #1
}
\newtcolorbox{cautionbox}[1][]{
  colback=red!3,
  colframe=red!50!black,
  boxrule=0.35mm,
  arc=1mm,
  left=2mm,
  right=2mm,
  top=1mm,
  bottom=1mm,
  #1
}
\newtcolorbox{examplebox}[1][]{
  colback=green!3,
  colframe=green!45!black,
  boxrule=0.35mm,
  arc=1mm,
  left=2mm,
  right=2mm,
  top=1mm,
  bottom=1mm,
  #1
}
\newtcolorbox{takeawaybox}[1][]{
  colback=VeryLightGray,
  colframe=RedTitle!55,
  boxrule=0.35mm,
  arc=1mm,
  left=2mm,
  right=2mm,
  top=1mm,
  bottom=1mm,
  #1
}
\newcommand{\compacttables}{\renewcommand{\arraystretch}{1.15}}

\setbeamertemplate{navigation symbols}{}
\setbeamertemplate{footline}{%
  \hfill%
  \usebeamercolor[fg]{page number in head/foot}%
  \usebeamerfont{page number in head/foot}%
  \insertframenumber\,/\,\inserttotalframenumber\kern1em\vskip2pt%
}

\makeatother

\author{Zhigang Feng}
% "date with time": compile timestamp so each build is identifiable.
\date{\today\ \DTMcurrenttime}


% Suppress redundant auto-generated metropolis section page; \sectiondivider frames are the dividers we keep.
\metroset{sectionpage=none}

\graphicspath{{./pic/}{./figures/}{../}}

\usetikzlibrary{positioning,arrows.meta,shapes,calc,fit,backgrounds}

%% Lec09 shared MAP slide, v2 (September 2026 reorganization): five parts,
%% eleven decks.  v1 (the July "five acts" map) is archived as
%% _archive/five_acts_2026-07/lec09_map_v1.tex.
%%
%% Input this file AFTER ../shared_preamble.tex (needs RedTitle, VeryLightGray,
%% DarkText) and call \LecNineMap{part}{sub} inside a frame:
%%   part in {1..5} highlights that part ("you are here") and, when sub names a
%%   sub-deck letter, that deck's chip -- \LecNineMap{2}{b} is T2b, \LecNineMap{4}{}
%%   is T4;  part = 0 renders the closing reprise with every part complete (the
%%   "Your Job Description" frame in T5d).
%% Each part carries its student question and its economic twin (the delegation-
%% economics thread of Lec09_Reorg_Plan_Review_2026-09-12.md, A3).
%% Filename deliberately does not match the build glob Lec*_T*.tex, so the build
%% script never tries to compile it standalone.

% Highlight chip #1 when the requested deck key equals #2 (e.g. 2b).
\newcommand{\lecmapchip}[2]{%
  \edef\lecmaptmp{#2}%
  \ifx\lecmapkey\lecmaptmp\tikzset{#1/.style={lecchipcur}}\fi}

\newcommand{\LecNineMap}[2]{%
% TikZ option lists are comma-split before TeX conditionals run, so every
% \ifnum/\ifx is resolved here, into named styles, before the picture starts.
\edef\lecmapkey{#1#2}%
\tikzset{lecparta/.style={}, lecpartb/.style={}, lecpartc/.style={},
         lecpartd/.style={}, lecparte/.style={},
         chipTone/.style={}, chipTtwoa/.style={}, chipTtwob/.style={},
         chipTthreea/.style={}, chipTthreeb/.style={}, chipTfour/.style={},
         chipTfivea/.style={}, chipTfiveb/.style={}, chipTfivec/.style={},
         chipTfived/.style={}}%
\ifnum#1=0
  \tikzset{lecparta/.style={lecpartdone}, lecpartb/.style={lecpartdone},
           lecpartc/.style={lecpartdone}, lecpartd/.style={lecpartdone},
           lecparte/.style={lecpartdone}, lecchip/.style={lecchipbase, lecchipdone}}%
\else
  \tikzset{lecchip/.style={lecchipbase}}%
  \ifnum#1=1 \tikzset{lecparta/.style={lecpartcur}}\fi
  \ifnum#1=2 \tikzset{lecpartb/.style={lecpartcur}}\fi
  \ifnum#1=3 \tikzset{lecpartc/.style={lecpartcur}}\fi
  \ifnum#1=4 \tikzset{lecpartd/.style={lecpartcur}}\fi
  \ifnum#1=5 \tikzset{lecparte/.style={lecpartcur}}\fi
  \lecmapchip{chipTone}{1}\lecmapchip{chipTtwoa}{2a}\lecmapchip{chipTtwob}{2b}%
  \lecmapchip{chipTthreea}{3a}\lecmapchip{chipTthreeb}{3b}\lecmapchip{chipTfour}{4}%
  \lecmapchip{chipTfivea}{5a}\lecmapchip{chipTfiveb}{5b}%
  \lecmapchip{chipTfivec}{5c}\lecmapchip{chipTfived}{5d}%
\fi
\begin{center}
\begin{tikzpicture}[
    part/.style={rectangle, rounded corners, draw=black!60, fill=VeryLightGray,
        minimum width=2.62cm, minimum height=0.72cm, text width=2.5cm,
        align=center, font=\scriptsize, text=DarkText},
    lecpartcur/.style={draw=RedTitle, very thick, fill=orange!20},
    lecpartdone/.style={draw=green!45!black, thick, fill=green!10},
    lecchipbase/.style={rectangle, rounded corners=1pt, draw=black!30, fill=white,
        inner xsep=1.5pt, inner ysep=1pt, font=\tiny, text=black!70},
    lecchipcur/.style={draw=RedTitle, fill=RedTitle!12, text=RedTitle, font=\tiny\bfseries},
    lecchipdone/.style={draw=green!45!black, fill=green!8},
    q/.style={font=\tiny\itshape, text width=2.7cm, align=center, text=black!75},
    twin/.style={font=\tiny, text width=2.7cm, align=center, text=blue!40!black},
    sat/.style={rectangle, rounded corners, draw=black!45, dashed, fill=white,
        text width=5.6cm, align=center, font=\tiny, text=black!70, minimum height=0.42cm},
    barlab/.style={font=\tiny\bfseries, text=black!60},
    arr/.style={-{Stealth[length=2mm]}, thick, gray!60}
]
    % --- five part boxes ---
    \node[part, lecparta] (p1) at (0,0)     {\textbf{9.1 SEE}};
    \node[part, lecpartb] (p2) at (2.85,0)  {\textbf{9.2 UNDER THE HOOD}};
    \node[part, lecpartc] (p3) at (5.7,0)   {\textbf{9.3 BUILD}};
    \node[part, lecpartd] (p4) at (8.55,0)  {\textbf{9.4 DESIGN \& VALIDATE}};
    \node[part, lecparte] (p5) at (11.4,0)  {\textbf{9.5 RUN \& GOVERN}};
    \draw[arr] (p1) -- (p2);
    \draw[arr] (p2) -- (p3);
    \draw[arr] (p3) -- (p4);
    \draw[arr] (p4) -- (p5);
    % --- sub-deck chips ---
    \node[lecchip, chipTone]    at (0,-0.6)     {T1};
    \node[lecchip, chipTtwoa]   at (2.55,-0.6)  {T2a};
    \node[lecchip, chipTtwob]   at (3.15,-0.6)  {T2b};
    \node[lecchip, chipTthreea] at (5.4,-0.6)   {T3a};
    \node[lecchip, chipTthreeb] at (6.0,-0.6)   {T3b};
    \node[lecchip, chipTfour]   at (8.55,-0.6)  {T4};
    \node[lecchip, chipTfivea]  at (10.5,-0.6)  {T5a};
    \node[lecchip, chipTfiveb]  at (11.1,-0.6)  {T5b};
    \node[lecchip, chipTfivec]  at (11.7,-0.6)  {T5c};
    \node[lecchip, chipTfived]  at (12.3,-0.6)  {T5d};
    % --- the student question each part answers ---
    \node[q] at (0,-1.08)     {what changed,\\ and why care?};
    \node[q] at (2.85,-1.08)  {how does it run?\\ how do I start?};
    \node[q] at (5.7,-1.08)   {how do I build one,\\ and call it?};
    \node[q] at (8.55,-1.08)  {can I design one\\ and prove it works?};
    \node[q] at (11.4,-1.08)  {run a project, catch\\ mistakes, stay safe};
    % --- its economic twin ---
    \node[twin] at (0,-1.62)     {generation got cheap;\\ verification is scarce};
    \node[twin] at (2.85,-1.62)  {what am I paying for?\\ tokens, scarce context};
    \node[twin] at (5.7,-1.62)   {dividing the labour:\\ specialists, boundaries};
    \node[twin] at (8.55,-1.62)  {write the contract,\\ audit the work};
    \node[twin] at (11.4,-1.62)  {hidden action: gates,\\ monitoring, priors};
    % --- "you are here" marker above the current part ---
    \ifnum#1=1 \node[font=\scriptsize\bfseries, text=RedTitle] at (0,0.62)
        {you are here\ $\blacktriangledown$};\fi
    \ifnum#1=2 \node[font=\scriptsize\bfseries, text=RedTitle] at (2.85,0.62)
        {you are here\ $\blacktriangledown$};\fi
    \ifnum#1=3 \node[font=\scriptsize\bfseries, text=RedTitle] at (5.7,0.62)
        {you are here\ $\blacktriangledown$};\fi
    \ifnum#1=4 \node[font=\scriptsize\bfseries, text=RedTitle] at (8.55,0.62)
        {you are here\ $\blacktriangledown$};\fi
    \ifnum#1=5 \node[font=\scriptsize\bfseries, text=RedTitle] at (11.4,0.62)
        {you are here\ $\blacktriangledown$};\fi
    \ifnum#1=0 \node[font=\scriptsize\bfseries, text=green!45!black] at (5.7,0.62)
        {all five parts complete};\fi
    % --- bottom band: the three questions of the lecture ---
    \fill[blue!10]   (-1.3,-2.37) rectangle (4.15,-2.07);
    \fill[green!10]  (4.4,-2.37)  rectangle (9.85,-2.07);
    \fill[orange!15] (10.1,-2.37) rectangle (12.7,-2.07);
    \node[barlab] at (1.425,-2.22)  {HOW IT WORKS};
    \node[barlab] at (7.125,-2.22)  {HOW TO BUILD \& USE IT};
    \node[barlab] at (11.4,-2.22)   {YOUR ROLE};
    % --- dashed satellites ---
    \node[sat] at (1.9,-2.8)
        {\textbf{Appendix TA} --- the dated landscape and setup details, read on demand};
    \node[sat] at (9.9,-2.8)
        {\textbf{Lab} Parts A--B, Design Lab scenarios $\to$ \textbf{Lec10} cases (same morning)};
\end{tikzpicture}
\end{center}
}


\title[AI for Econ Research]{AI for Economic Research: Dynamic Models, Language, and Agents\\
Lecture 9.2b: Your Setup --- Files, Commands, First Session}

\begin{document}

\makebeamertitle

% ---- NEW · TIER: READ
\begin{frame}{Lecture 9 in One Map: Your Setup}
\textbf{This deck puts the loop on your machine: install it, learn its files and commands, and run a first session.}

\LecNineMap{2}{b}

\vspace{0.1cm}
\footnotesize\textit{Route: terminal, install, and Git $\to$ the files an agent reads $\to$ the control surface $\to$ your first session $\to$ tokens and the context window.}
\end{frame}


%=============================================================================
\section{Terminal, Install, Git}
%===========================================================

\sectiondivider{Terminal, Install, Git}{The pieces, one install, and the safety net}

% ---- T2-F05 · KEEP · TIER: READ · from T2:124
% MiniCourse refinement: terminal-workflow vocabulary from M4_T4
\begin{frame}{Terminal Vocabulary in 60 Seconds}
\small
\begin{center}
\begin{tabular}{|p{2.7cm}|p{8.5cm}|}
\hline
\textbf{Term} & \textbf{Meaning} \\ \hline
\textbf{Terminal} & The text window where you type commands instead of clicking buttons. \\ \hline
\textbf{Shell} & The program that interprets commands. macOS often uses \texttt{zsh}; many examples are written in \texttt{bash}. \\ \hline
\textbf{Bash} & A common shell language. When Claude Code says it will use \texttt{Bash}, it means ``run this command in the terminal.'' \\ \hline
\textbf{Working directory} & The folder commands run from. \texttt{cd my-project} changes it; \texttt{pwd} prints it. \\ \hline
\textbf{PATH} & The list of folders your shell searches when you type a command such as \texttt{python} or \texttt{claude}. \\ \hline
\end{tabular}
\end{center}

\vspace{0.15cm}
\textbf{Why this matters:} terminal agents act through shell commands. If the working directory or \texttt{PATH} is wrong, the agent may be competent but operating in the wrong place.
\end{frame}

% ---- T2-F06 · KEEP · TIER: READ · from T2:142
\begin{frame}{How the Terminal Pieces Fit Together}
\textbf{One picture for the five words: you type into the terminal; the shell finds the program via PATH and runs it in the working directory.}

\vspace{0.15cm}

\begin{center}
\begin{tikzpicture}[
    tb/.style={rectangle, rounded corners, draw=black!60, fill=VeryLightGray,
        minimum width=2.6cm, minimum height=1.05cm, text width=2.5cm,
        align=center, font=\scriptsize},
    side/.style={rectangle, rounded corners, draw=RedTitle!60, fill=blue!5,
        minimum width=3.4cm, minimum height=0.85cm, text width=3.3cm,
        align=center, font=\scriptsize},
    arr/.style={-{Stealth[length=2mm]}, thick, gray!60},
    lab/.style={font=\tiny\ttfamily, text=black!60}
]
    \node[tb] (you) at (0,0) {\textbf{You}\\ type a command:\\ \texttt{claude}};
    \node[tb] (term) at (3.3,0) {\textbf{Terminal}\\ the window that carries text in and out};
    \node[tb] (shell) at (6.6,0) {\textbf{Shell} (\texttt{zsh}/\texttt{bash})\\ interprets the line};
    \node[tb] (prog) at (9.9,0) {\textbf{Program runs}\\ \texttt{claude}, \texttt{python}, \texttt{git}};
    \node[side] (path) at (6.6,1.35) {\textbf{PATH:} the list of folders searched to \emph{find} the program};
    \node[side] (wd) at (9.9,-1.35) {\textbf{Working directory:} the folder the program runs \emph{in} (reads/writes here)};
    \draw[arr] (you) -- (term);
    \draw[arr] (term) -- (shell);
    \draw[arr] (shell) -- node[lab, above] {found it} (prog);
    \draw[arr, dashed] (path) -- (shell);
    \draw[arr, dashed] (wd) -- (prog);
\end{tikzpicture}
\end{center}

\begin{itemize}
    \item When you type \texttt{claude}, the shell finds it via \texttt{PATH} and starts it \emph{in your current folder} --- that is why \texttt{cd my-project} comes first
    \item The agent's \texttt{Bash} tool is this same picture one level down: the agent typing into the same shell, in the same working directory
\end{itemize}
\end{frame}

% ---- T2-F07 · EDIT · TIER: LAB · from T2:185
% TODO(Phase B): Review A2: the three durable steps shown for each tested harness (Claude Code, Gemini CLI, Codex CLI), commands from CLI_FACTS_2026-09.md.
\begin{frame}[fragile]{Install Once: Three Durable Steps}
\textbf{Setup has exactly three steps, and only the first command ever changes.}

\vspace{0.2cm}

\begin{enumerate}
    \item \textbf{Install the CLI} --- one installer command for your OS (native installer script; Homebrew and \texttt{npm} also work):
\begin{lstlisting}[style=pythonstyle,language={},basicstyle=\ttfamily\scriptsize,numbers=none]
curl -fsSL https://claude.ai/install.sh | bash
\end{lstlisting}
    \item \textbf{Authenticate once} --- run \texttt{claude}; a browser window opens; log in with your account or API key
    \item \textbf{Work from the project} --- every session starts the same way:
\begin{lstlisting}[style=pythonstyle,language={},basicstyle=\ttfamily\scriptsize,numbers=none]
cd my-project
claude
\end{lstlisting}
\end{enumerate}

\vspace{0.15cm}

\begin{takeawaybox}
\footnotesize
\textbf{Install once per machine. Authenticate once per account. Initialize every session per project --- automatically (next frames).}
\end{takeawaybox}

\vspace{0.1cm}
\footnotesize\textit{Durable vs.\ dated: this 3-step shape survives every product update; exact commands, subscription tiers, and the Windows/OneDrive caveats are dated detail $\to$ Appendix TA and the official docs.}
\end{frame}

% ---- T1-F14 · MERGE-PENDING(T1-F14+T1-F15) · TIER: CORE · from T1:325
% TODO(Phase B): One frame: the time-machine TikZ from T1-F14 above T1-F15's risk/benefit columns.
\begin{frame}{What Git Is: A Time Machine for Your Project}
\textbf{Git is a version-control system: it keeps snapshots of a project folder so you can always see what changed and go back.}

\vspace{0.15cm}

\begin{center}
\begin{tikzpicture}[
    gb/.style={rectangle, rounded corners, draw=black!60, fill=VeryLightGray,
        minimum width=2.75cm, minimum height=1.15cm, text width=2.65cm,
        align=center, font=\scriptsize},
    cmd/.style={font=\tiny\ttfamily, text=RedTitle},
    arr/.style={-{Stealth[length=2mm]}, thick, gray!60}
]
    \node[gb] (wd) at (0,0) {\textbf{Working directory}\\ your files, edited freely (by you or the agent)};
    \node[gb] (st) at (3.9,0) {\textbf{Staging area}\\ what the next snapshot will include};
    \node[gb] (hist) at (7.8,0) {\textbf{History}\\ snapshot 1 $\leftarrow$ 2 $\leftarrow$ 3 $\cdots$\\ \tiny lives in a hidden \texttt{.git} folder};
    \node[gb, dashed, fill=white] (rem) at (11.7,0) {\textbf{Remote (GitHub)}\\ optional copy: backup, sharing};
    \draw[arr] (wd) -- node[above, cmd] {git add} (st);
    \draw[arr] (st) -- node[above, cmd] {git commit} (hist);
    \draw[arr, dashed] (hist) -- node[above, cmd] {git push} (rem);
\end{tikzpicture}
\end{center}

\vspace{0.2cm}

\begin{itemize}
    \item A \textbf{commit} is a snapshot of the whole project at a moment you chose; the history is a chain of snapshots you can list (\texttt{git log}), compare (\texttt{git diff}), and return to
    \item Everything is local, plain files, no server required --- the \textbf{remote} is optional
    \item \textbf{Why it fits agent work:} the agent edits fast and wide; snapshots give you a clean before/after boundary around every session
\end{itemize}
\end{frame}

% ---- T1-F15 · MERGE-PENDING(T1-F14+T1-F15) · TIER: CORE · from T1:357
\begin{frame}{Git Is the Safety Net}

\begin{columns}[T]
\begin{column}{0.48\textwidth}
\textbf{The risk without Git:}
\vspace{0.15cm}
\begin{itemize}
    \item Claude Code can edit dozens of files autonomously in one session
    \item OneDrive and Dropbox sync by \textit{copying}---two simultaneous writes create corrupt merge files
    \item No undo: if the agent makes bad changes, there is no ``revert'' without Git
    \item No audit trail: you cannot see what changed or when
\end{itemize}
\end{column}
\begin{column}{0.48\textwidth}
\textbf{What Git gives you:}
\vspace{0.15cm}
\begin{itemize}
    \item \texttt{git diff} shows exactly what the agent changed before you accept it
    \item \texttt{git checkout -{-} .} reverts all uncommitted changes instantly---your emergency brake
    \item \texttt{git log} shows every change with a timestamp and message
    \item Works offline; cloud sync races with the agent's rapid file writes
\end{itemize}
\end{column}
\end{columns}

\vspace{0.25cm}

\textbf{Practical rule:} keep large data files on OneDrive; keep code and scripts in Git. Mixing them is workable if you are careful --- but do not let OneDrive sync a \emph{live repository} (the hidden \texttt{.git} folder): the two sync models collide and can corrupt the history. Keep the repo outside the synced folder, or exclude \texttt{.git} from sync.
\end{frame}

% ---- T1-F16 · KEEP · TIER: LAB · from T1:387
\begin{frame}[fragile]{Git Setup: A Four-Step Roadmap (Hand It to Your Agent)}
\textbf{You do not need to memorize commands. You need the roadmap --- any LLM can walk you through each step.}

\vspace{0.1cm}

\begin{columns}[T]
\begin{column}{0.48\textwidth}
\textbf{Step 0 --- check the install:}
\begin{lstlisting}[style=pythonstyle,language={},basicstyle=\ttfamily\scriptsize,numbers=none]
git --version
\end{lstlisting}
{\footnotesize (if missing: git-scm.com or your package manager)}

\vspace{0.15cm}

\textbf{Step 1 --- identify yourself} (once per machine):
\begin{lstlisting}[style=pythonstyle,language={},basicstyle=\ttfamily\scriptsize,numbers=none]
git config --global user.name "Your Name"
git config --global user.email "you@edu.cn"
\end{lstlisting}
\end{column}
\begin{column}{0.48\textwidth}
\textbf{Step 2 --- start the project} (once per project):
\begin{lstlisting}[style=pythonstyle,language={},basicstyle=\ttfamily\scriptsize,numbers=none]
cd my-project
git init
echo "data/ output/ *.pdf" > .gitignore
git add . && git commit -m "initial setup"
\end{lstlisting}

\vspace{0.15cm}

\textbf{Step 3 --- the session rhythm} (every session):
\begin{lstlisting}[style=pythonstyle,language={},basicstyle=\ttfamily\scriptsize,numbers=none]
git add -A && git commit -m "pre-session"
# ... agent works, you validate ...
git add -A && git commit -m "validated: ..."
\end{lstlisting}
\end{column}
\end{columns}

\vspace{0.15cm}

\begin{takeawaybox}
\footnotesize
\textbf{Emergency brake:} \texttt{git checkout -{-} .} undoes \emph{all} uncommitted changes instantly. And note: this setup is itself a perfect first agent task --- paste the roadmap into the agent and say \emph{``walk me through steps 0--3 for this folder.''}
\end{takeawaybox}
\end{frame}


%=============================================================================
\section{Files: Where the Agent Reads From}
%===========================================================

\sectiondivider{Files: Where the Agent Reads From}{Global versus project, and what loads when you start}

% ---- T2-F08 · EDIT · TIER: READ · from T2:214
\begin{frame}{Files First: Global vs.\ Local}

Two layers of configuration. \textbf{Local overrides global.}

\vspace{0.2cm}

\begin{center}
\footnotesize
\begin{tabular}{|l|l|p{4.9cm}|}
\hline
\textbf{Path (Claude Code example)} & \textbf{Scope} & \textbf{What lives here} \\ \hline
\texttt{\textasciitilde/.claude/settings.json} & Global & Default permissions, model preference \\ \hline
\texttt{\textasciitilde/.claude/CLAUDE.md} & Global & Your cross-project preferences \\ \hline
\texttt{\textasciitilde/.claude/skills/} & Global & Skills available in every project \\ \hline
\texttt{\textasciitilde/.claude/projects/\ldots/memory/} & Global & Auto memory: the tool's own per-project notes \\ \hline
\hline
\texttt{<project>/CLAUDE.md} & Local & Project brain, read every session \\ \hline
\texttt{<project>/.claude/settings.json} & Local & Project permissions (overrides global) \\ \hline
\texttt{<project>/.claude/skills/} & Local & This project's slash commands \\ \hline
\texttt{<project>/.claude/agents/} & Local & Specialized reviewer roles \\ \hline
\texttt{<project>/.claude/rules/} & Local & Always-on guardrails \\ \hline
\end{tabular}
\end{center}

\vspace{0.15cm}

\footnotesize\textbf{Other agents follow analogous layouts.} Codex CLI uses \texttt{AGENTS.md}; Gemini CLI uses \texttt{GEMINI.md}. The pattern --- global defaults plus a project-specific system prompt --- is durable across products.

\vspace{0.05cm}

\textbf{Rule:} global dotfolder $=$ your personal defaults (apply everywhere). Local dotfolder $=$ the project's rules (apply only here). T5b fills the local layer with research content.
\end{frame}

% ---- NEW · TIER: CORE
% TODO(Phase B): Review A2.5 (replaces v1's .claude-only tree and TA's mapping table): a verified table from CLI_FACTS_2026-09.md -- instructions file (CLAUDE.md / GEMINI.md / AGENTS.md), skills, custom sub-agents (or "not native -- use rung 3"), settings/policy, memory -- per tested harness, with version and date.
% TODO(Phase B): Live examples: ROOT/.claude/agents/, labs/Lec09_Agent_Lab/.claude/agents/, demos/M4_olg_5step_demo/.claude/skills/olg-5step/SKILL.md. Draw only paths that exist.
\begin{frame}{Project Files Across Harnesses}
\textbf{Placeholder --- written in Phase B.} The specification is in the source comment above this frame.
\end{frame}

% ---- T2-F09 · KEEP · TIER: READ · from T2:247
\begin{frame}{Initialization: What Happens When You Type \texttt{claude}}

Five steps run \textbf{before you type your first message} --- each one loads a file from the previous frame:

\vspace{0.2cm}

\begin{enumerate}
    \item \textbf{Read global config} --- \texttt{\textasciitilde/.claude/settings.json}: permission defaults, model preference

    \vspace{0.1cm}

    \item \textbf{Detect project context} --- searches current directory (and parents) for \texttt{CLAUDE.md} and \texttt{.claude/}

    \vspace{0.1cm}

    \item \textbf{Load project memory} --- reads \texttt{CLAUDE.md} (and the tool's auto memory for this project) into the context window

    \vspace{0.1cm}

    \item \textbf{Auto-load all rules} --- every \texttt{.claude/rules/*.md} file is loaded silently

    \vspace{0.1cm}

    \item \textbf{Ready} --- context window pre-populated before you say anything
\end{enumerate}

\vspace{0.2cm}

\textbf{Analogy:} steps 3--4 are a professor re-reading lecture notes before class. Students arrive; context is already loaded.

\vspace{0.1cm}

\textbf{Cost warning:} steps 3--4 consume tokens before you type anything. Keep \texttt{CLAUDE.md} and rules files concise.
\end{frame}

% ---- T4-F10 · MOVE · TIER: READ · from T4:271
\begin{frame}{Three Memories, Three Writers}
\textbf{Three different persistence ideas should not be confused.}

\vspace{0.2cm}

\begin{center}
\small
\begin{tabular}{|p{2.4cm}|p{2.4cm}|p{5.3cm}|}
\hline
\textbf{Mechanism} & \textbf{Who writes it} & \textbf{Why it exists} \\ \hline
\texttt{CLAUDE.md} & you / team & durable project instructions and conventions \\ \hline
Auto memory & the tool & per-project learned patterns and corrections \\ \hline
Optional course log & you / course & explicit progress notes such as \texttt{progress.md} \\ \hline
\end{tabular}
\end{center}

\vspace{0.25cm}

\textbf{Course recommendation:} rely on \texttt{CLAUDE.md} for shared rules, treat auto memory as helpful but not authoritative, and keep explicit logs when you want students to inspect the evolution of a project.
\end{frame}

% ---- T2c-F12 · EDIT · TIER: READ · from T2c:425
% TODO(Phase B): Review A2: the table rows are Claude Code objects; add one line on how the other tested harnesses scope the same thing (from CLI_FACTS).
\begin{frame}{Where Those Four Agents Lived}
\textbf{Three different things get called ``the agent'', and they have three different lifetimes. Only one of them is yours.}

\vspace{0.1cm}

\begin{center}
\footnotesize
\renewcommand{\arraystretch}{1.24}
\begin{tabular}{|p{2.8cm}|p{2.3cm}|p{3.5cm}|p{1.8cm}|}
\hline
 & \textbf{Written by} & \textbf{Lives in} & \textbf{Survives?} \\ \hline
the \emph{type} (\texttt{Explore}, \texttt{Plan}) & the vendor, built in & the harness, delivered as a list in the system prompt & yes, but not yours to edit \\ \hline
the \emph{brief} & \textbf{the model}, in the tool call & nowhere --- four tool-call arguments & \textbf{no} \\ \hline
the \emph{report} & the sub-agent & the message list, plus session scratch & no \\ \hline
\texttt{code-reviewer}, \texttt{math-agent}, \ldots & \textbf{you} & \texttt{.claude/agents/} in the lab & yes, versioned \\ \hline
\end{tabular}
\end{center}

\vspace{0.1cm}

\begin{cautionbox}
\footnotesize\textbf{Where you put the folder decides whether the specialist exists.} \texttt{.claude/agents/} resolves relative to the directory the session starts \emph{in} --- the global-vs-local rule above, applied to agents.
\end{cautionbox}
\end{frame}


%=============================================================================
\section{The Control Surface}
%===========================================================

\sectiondivider{The Control Surface}{Commands, modes, keys, and permissions}

% ---- T2-F10 · KEEP · TIER: READ · from T2:282
\begin{frame}{The CLI at a Glance: One Screen, Four Zones}
\textbf{Everything you will ever look at in a terminal-agent session is one screen with four zones.}

\vspace{0.1cm}

\begin{center}
\begin{tikzpicture}[
    zone/.style={rectangle, draw=black!50, align=left, font=\ttfamily\tiny, inner sep=3pt},
    lab/.style={font=\scriptsize, text=RedTitle, align=left},
    larr/.style={-{Stealth[length=1.6mm]}, thin, RedTitle!70}
]
    % terminal window
    \node[zone, fill=white, minimum width=8.6cm, minimum height=2.15cm, text width=8.3cm] (tr) at (0,0)
        {You> Parse the roster page into a seed CSV\\[1pt]
         claude> {[Read sources/fellows\_page.html]}\\
         \phantom{claude>} {[Write data/fellows\_seed.csv]}\\
         \phantom{claude>} Parsed 882 rows; flagged 2 unparseable\\
         \phantom{claude>} election years in notes/batch\_log.md};
    \node[zone, fill=orange!10, draw=orange!70, minimum width=8.6cm, minimum height=0.55cm, text width=8.3cm, below=0.06cm of tr] (perm)
        {Allow Bash(python scripts/parse\_roster.py)?\ \ [y]es / [n]o / always};
    \node[zone, fill=blue!4, minimum width=8.6cm, minimum height=0.5cm, text width=8.3cm, below=0.06cm of perm] (inp)
        {> \_};
    \node[zone, fill=VeryLightGray, minimum width=8.6cm, minimum height=0.45cm, text width=8.3cm, below=0.06cm of inp] (stat)
        {mode: plan\ \ $\cdot$\ \ context used: 34\%\ \ $\cdot$\ \ session cost so far: /cost};
    % labels
    \node[lab, right=0.45cm of tr.east, text width=3.2cm] (l1)
        {\textbf{1 Transcript}\\ \scriptsize dialogue + bracketed tool traces (the JSON calls, printed)};
    \node[lab, right=0.45cm of perm.east, text width=3.2cm] (l2)
        {\textbf{2 Permission prompt}\\ \scriptsize approve actions one by one};
    \node[lab, right=0.45cm of inp.east, text width=3.2cm] (l3)
        {\textbf{3 Your input}\\ \scriptsize questions, commands, /slash};
    \node[lab, right=0.45cm of stat.east, text width=3.2cm] (l4)
        {\textbf{4 Status line}\\ \scriptsize mode, context usage, cost};
    \draw[larr] (l1.west) -- (tr.east);
    \draw[larr] (l2.west) -- (perm.east);
    \draw[larr] (l3.west) -- (inp.east);
    \draw[larr] (l4.west) -- (stat.east);
\end{tikzpicture}
\end{center}

\vspace{0.1cm}
\footnotesize\textit{Zone 1's \texttt{[Read ...]} traces are the JSON tool calls of the earlier frame, printed for you --- the audit trail is on screen by default. Layout details differ across products and versions; the four zones are durable.}
\end{frame}

% ---- T2-F11 · EDIT · TIER: READ · from T2:326
\begin{frame}{Commands That Earn Their Keep (1): Framing and Steering}
\textbf{These six decide what the agent attempts --- and how hard it thinks.}

\vspace{0.15cm}

\begin{center}
\scriptsize
\renewcommand{\arraystretch}{1.05}
\begin{tabular}{|l|p{9.4cm}|}
\hline
\textbf{Command} & \textbf{What it does --- and why it earns a place} \\ \hline
\multicolumn{2}{|l|}{\textbf{Before the work starts}} \\ \hline
\texttt{/init} & Writes a first \texttt{CLAUDE.md} from the repo --- the brief that is re-read every session (T5b). Generating it beats writing it cold \\ \hline
\texttt{/plan} & Explores and drafts; executes \textbf{nothing} until you approve --- so a wrong turn costs a re-read, not a \texttt{git reset} \\ \hline
\multicolumn{2}{|l|}{\textbf{How hard it thinks}} \\ \hline
\texttt{/model} & Switches model mid-session, context preserved. A derivation and a variable rename do not deserve the same engine \\ \hline
\texttt{/effort} & Five levels, \texttt{low} $\to$ \texttt{max}: the biggest quality--cost dial you hold. Raise it for a proof or a subtle bug; lower it for mechanical edits \\ \hline
\texttt{/fast} & Toggles fast mode: same model, faster output --- \textbf{not} a downgrade to a smaller one. For when latency, not depth, is the bottleneck \\ \hline
\multicolumn{2}{|l|}{\textbf{How you get words in}} \\ \hline
\texttt{/voice} & Dictation (\texttt{hold} / \texttt{tap} / \texttt{off}). You specify a task far faster than you type it --- and thin prompts are where these tools underperform \\ \hline
\end{tabular}
\end{center}

\vspace{0.1cm}

\footnotesize\textbf{Spend effort on \textit{decisions}, save it on \textit{execution}:} max effort on a long context can cost an order of magnitude more than the default.
\end{frame}

% ---- T2-F12 · EDIT · TIER: READ · from T2:354
% TODO(Phase B): Add one pointer row: "/agents -- inspect which specialists are in scope (T3b)". T2b owns the twelve commands; T3b owns the rung surfaces.
\begin{frame}{Commands That Earn Their Keep (2): Context, Cost, and Undo}
\textbf{The other six manage your two budgets --- the window and the meter --- and undo mistakes.}

\vspace{0.15cm}

\begin{center}
\scriptsize
\renewcommand{\arraystretch}{1.05}
\begin{tabular}{|l|p{9.4cm}|}
\hline
\textbf{Command} & \textbf{What it does --- and why it earns a place} \\ \hline
\multicolumn{2}{|l|}{\textbf{Mind the window}} \\ \hline
\texttt{/context} & Draws current context use as a colored grid. Makes the invisible constraint visible: you see what is eating the window \textit{before} quality degrades \\ \hline
\texttt{/compact} & Summarizes the conversation to free space. Keeps the thread, drops the detail --- one you choose beats one that fires mid-derivation \\ \hline
\texttt{/clear} & New session, empty context; the old one stays on disk. The right move \textit{between} unrelated tasks, which otherwise contaminate each other \\ \hline
\multicolumn{2}{|l|}{\textbf{Mind the meter --- and the undo button}} \\ \hline
\texttt{/usage} & Session cost, plan usage, activity (alias \texttt{/cost}). Turns an invisible meter into a number: check \textit{before} a long agentic run, not after \\ \hline
\texttt{/rewind} & Restores files to an earlier checkpoint --- the safety net that makes accept-edits mode reasonable \\ \hline
\texttt{/resume} & Reopens an earlier conversation. Research is not one sitting; yesterday's context is an asset, not something to re-explain \\ \hline
\end{tabular}
\end{center}

\vspace{0.02cm}

\begin{takeawaybox}
\scriptsize
\textbf{Four dials, not twelve commands.} What it reads (\texttt{/context}, \texttt{/compact}, \texttt{/clear}) $\cdot$ how hard it thinks (\texttt{/model}, \texttt{/effort}, \texttt{/fast}) $\cdot$ how much it may do unasked (modes, \texttt{/permissions}) $\cdot$ what it costs (\texttt{/usage}).
\end{takeawaybox}
\end{frame}

% ---- T2-F13 · EDIT · TIER: READ · from T2:384
\begin{frame}{Modes, Keys, and the Rest of the Surface}
\textbf{Commands are only one of three ways to steer. The other two are faster.}

\vspace{0.15cm}

\begin{columns}[T]
\begin{column}{0.49\textwidth}
\textbf{Modes} --- cycle with \texttt{Shift+Tab}:

\vspace{0.1cm}
{\scriptsize
\begin{tabular}{|l|p{3.8cm}|}
\hline
Plan & Read-only: explores and drafts, executes \textbf{nothing} until you approve. Start every non-trivial task here \\ \hline
Accept-edits & Auto-approves routine file edits. Only after trust is earned \\ \hline
Default & Asks before each action \\ \hline
\end{tabular}
}

\vspace{0.2cm}
\textbf{Keys:}

\vspace{0.1cm}
{\scriptsize
\begin{tabular}{|l|p{3.8cm}|}
\hline
\texttt{Esc} & Interrupt the agent mid-task, immediately \\ \hline
\texttt{Esc Esc} & Rewind to an earlier checkpoint \\ \hline
\texttt{!} & Run a shell command yourself; the output joins the context \\ \hline
\texttt{@} & Point at a file by path \\ \hline
\texttt{Ctrl+R} & Search your prompt history \\ \hline
\end{tabular}
}
\end{column}
\begin{column}{0.49\textwidth}
\textbf{Worth knowing, reached for less often:}

\vspace{0.1cm}
{\scriptsize
\begin{tabular}{|l|p{3.7cm}|}
\hline
\texttt{/memory} & Edit \texttt{CLAUDE.md} and memory settings directly \\ \hline
\texttt{/permissions} & Manage allow and deny rules --- next frame \\ \hline
\texttt{/status} & Version, model, account, connectivity \\ \hline
\texttt{/export} & Save the conversation to a file: an audit trail of how a result was produced \\ \hline
\texttt{/hooks} & Inspect what runs automatically on tool events (T5b) \\ \hline
\texttt{/doctor} & Health-check the install when something is off \\ \hline
\end{tabular}
}

\vspace{0.15cm}
\footnotesize \texttt{/help} lists every command; typing \texttt{/} alone opens the menu.
\end{column}
\end{columns}

\vspace{0.1cm}
\footnotesize\textit{Durable vs.\ dated: the four dials --- context, effort, permission, cost --- outlive any release. Exact command names, effort levels, and model aliases are dated detail $\to$ \texttt{/help} and the official docs.}
\end{frame}

% ---- T2-F14 · KEEP · TIER: CORE · from T2:443
\begin{frame}[fragile]{Permissions: The Blast-Radius Dial}

\begin{columns}[T]
\begin{column}{0.44\textwidth}
\textbf{The agent's toolbox:}

\vspace{0.15cm}
\footnotesize
\begin{tabular}{|l|p{2.8cm}|}
\hline
\textbf{Tool} & \textbf{What it does} \\ \hline
\texttt{Bash} & Run any shell command \\ \hline
\texttt{Read}/\texttt{Grep} & Read and search files \\ \hline
\texttt{Write} & Create or overwrite a file \\ \hline
\texttt{Agent} & Spawn a child instance \\ \hline
\end{tabular}

\vspace{0.15cm}
\textbf{``Which Python ran this?''}
\begin{itemize}
    \item \scriptsize the agent's \texttt{Bash} inherits your shell's \texttt{\$PATH}: if you ran \texttt{conda activate myenv} before \texttt{claude}, \textit{that} Python runs
    \item \scriptsize add a \texttt{which python} check to \texttt{CLAUDE.md}
\end{itemize}
\end{column}
\begin{column}{0.53\textwidth}
\textbf{The permission model:}

\vspace{0.1cm}
\footnotesize
\begin{itemize}
    \item \textbf{Default: ask.} You see the exact command before it runs
    \item \textbf{Allowlist} trusted patterns in \texttt{.claude/settings.json}:
\end{itemize}
\begin{lstlisting}[style=pythonstyle,language={},basicstyle=\ttfamily\tiny,numbers=none]
"permissions": {
  "allow": ["Bash(python *)",
            "Bash(git *)"],
  "deny":  ["Read(.env)",
            "Bash(rm -rf *)"] }
\end{lstlisting}
\begin{itemize}
    \item \textbf{Bypass-permissions} skips all prompts: isolated sandboxes only
\end{itemize}
\end{column}
\end{columns}

\vspace{0.08cm}
\begin{conceptbox}
\footnotesize
\textbf{Think of authorizing a new RA.} Default $=$ every action needs your signature; allowlist $=$ standing approval for routine work; deny $=$ never allowed; bypass $=$ handing over your own login. \emph{``Blast radius''} $=$ how much damage one wrong action could do before you notice.
\end{conceptbox}
\end{frame}

% ---- NEW · TIER: READ
% TODO(Phase B): Optional (v1 T2b item 17): the commands, modes, keys and permission settings on one page, one column per tested harness, linked to the official docs; every entry from CLI_FACTS_2026-09.md.
\begin{frame}{One-Page Cheat Sheet}
\textbf{Placeholder --- written in Phase B.} The specification is in the source comment above this frame.
\end{frame}


%=============================================================================
\section{Your First Session}
%===========================================================

\sectiondivider{Your First Session}{A real walkthrough on the fellows roster}

% ---- T2-F15 · KEEP · TIER: LAB · from T2:503
\begin{frame}[fragile]{First Session (1): Launch on a Real Project}
\textbf{Initialization in the abstract $\to$ what it actually looks like, on the fellows project from T1.}

\vspace{0.08cm}
\begin{lstlisting}[basicstyle=\ttfamily\scriptsize,numbers=none]
$ cd ~/research/es-fellows
$ claude
[reads ~/.claude/settings.json]
[reads ./CLAUDE.md]
[loads .claude/rules/provenance.md]
[ready]

You> Parse sources/fellows_page.html into data/fellows_seed.csv
     with columns: name, affiliation, election_year, profile_url.

claude> [Read sources/fellows_page.html]
        [Write data/fellows_seed.csv]
        Parsed 882 rows. Two election years did not parse
        (flagged in notes/batch_log.md). Spot-check 5 rows?
\end{lstlisting}

\vspace{0.05cm}
\footnotesize
\begin{enumerate}
    \item Initialization is automatic; no flags needed for the first session
    \item The agent reads context (rules, references) \emph{before} acting --- visible as bracketed \texttt{[Read ...]} traces
    \item Each trace is a JSON tool call that the harness translated into a real disk operation
    \item Note the agent's last line: a competent session \emph{ends turns with a checkable claim}, not just ``done''
\end{enumerate}
\end{frame}

% ---- T2-F16 · EDIT · TIER: LAB · from T2:534
\begin{frame}{First Session (2): Pen and Paper, Then Plan Mode}

\begin{columns}[T]
\begin{column}{0.52\textwidth}
\small
\begin{enumerate}
    \item[\textbf{0.}] \textbf{Before the agent---pen~+~paper}\\
    Write one sentence: \textit{``I want to compute X using method Y for question Z.''}\\
    {\footnotesize Fellows version: ``a provenance-aware seed roster of all current ES Fellows, from the official page only.''}

    \vspace{0.1cm}

    \item[\textbf{1.}] \textbf{Enter plan mode (\texttt{Shift+Tab})}\\
    Describe your goal. The agent drafts a plan. \textit{Nothing runs yet.}

    \vspace{0.1cm}

    \item[\textbf{2.}] \textbf{Review and approve}\\
    Read the plan. Ask questions. Revise. Then approve it.
\end{enumerate}
\end{column}
\begin{column}{0.44\textwidth}
\begin{tcolorbox}[colback=blue!3, colframe=RedTitle!60]
\footnotesize
\textbf{Why this matters:} plan review is the cheapest place to catch a wrong task definition.
\end{tcolorbox}
\vspace{0.2cm}
\footnotesize
\textbf{Common beginner mistake:} skipping Step~0. Without a clear goal, the agent will confidently build the wrong thing.

\vspace{0.15cm}
\textit{What exactly to check before approving a plan: T5a gives the five plan-grading criteria.}
\end{column}
\end{columns}
\end{frame}

% ---- T2-F17 · KEEP · TIER: LAB · from T2:570
\begin{frame}{First Session (3): Implement, Review, Iterate}
\begin{columns}[T]
\begin{column}{0.52\textwidth}
\small
\begin{enumerate}
    \item[\textbf{3.}] \textbf{Implement}\\
    The agent executes the approved plan. Press \texttt{Esc} to stop if wrong.

    \vspace{0.1cm}

    \item[\textbf{4.}] \textbf{Review output}\\
    Does it make economic sense? Re-derive one number by hand.\\
    {\footnotesize Fellows version: open 5 random rows; follow each source URL.}

    \vspace{0.1cm}

    \item[\textbf{5.}] \textbf{Iterate}\\
    Fix, extend one feature, compact context, and return to step~1.
\end{enumerate}
\end{column}
\begin{column}{0.44\textwidth}
\begin{center}
\begin{tikzpicture}[
    stagebox/.style={rectangle, draw=black, fill=VeryLightGray, rounded corners,
                 minimum width=2.8cm, minimum height=0.65cm,
                 text width=2.6cm, align=center, font=\footnotesize},
    arrow/.style={-{Stealth[length=2mm]}, thick}
]
    \node[stagebox] (s0) {0.\ Pen + paper};
    \node[stagebox, below=0.22cm of s0] (s1) {1.\ Plan mode};
    \node[stagebox, below=0.22cm of s1] (s2) {2.\ Approve};
    \node[stagebox, below=0.22cm of s2] (s3) {3.\ Implement};
    \node[stagebox, below=0.22cm of s3] (s4) {4.\ Review};
    \draw[arrow] (s0) -- (s1);
    \draw[arrow] (s1) -- (s2);
    \draw[arrow] (s2) -- (s3);
    \draw[arrow] (s3) -- (s4);
    \draw[arrow] (s4.east) -- ++(0.6,0) -- ++(0,2.95) -- (s1.east);
    \node[right=0.4cm of s3, font=\tiny, rotate=90] {iterate};
\end{tikzpicture}
\end{center}
\vspace{0.2cm}
\footnotesize
\textbf{Practical takeaway:} after each implementation pass, force one human check before you ask for the next extension.
\end{column}
\end{columns}
\end{frame}

% ---- T2-F18 · KEEP · TIER: READ · from T2:618
\begin{frame}{What the First Session Leaves Behind}
\textbf{Close the terminal. What survives is not the conversation --- it is the project state on disk.}

\vspace{0.2cm}

\begin{columns}[T]
\begin{column}{0.5\textwidth}
\textbf{In the folder now:}
\footnotesize
\begin{itemize}
    \item \texttt{data/fellows\_seed.csv} --- the artifact (882 rows)
    \item \texttt{notes/batch\_log.md} --- what parsed, what was flagged
    \item \texttt{CLAUDE.md} --- two new lines: the parse rules you settled on
    \item \texttt{.claude/rules/provenance.md} --- the source rule, now always-on
\end{itemize}
\end{column}
\begin{column}{0.46\textwidth}
\textbf{Why this is the whole game:}
\footnotesize
\begin{itemize}
    \item chat history is session-bound (resumable, but do not rely on it); \textbf{files persist and compose}
    \item next session, initialization reloads all of this \emph{automatically} --- batch 2 starts where batch 1 ended
    \item ``the project files hold the rules and state across sessions'' --- the public-web analog of \texttt{CLAUDE.md} (Lec10 T4)
\end{itemize}
\end{column}
\end{columns}

\vspace{0.25cm}

\begin{takeawaybox}
\footnotesize
\textbf{A session is good when it improves the project state, not when the conversation felt smart.}
\end{takeawaybox}
\end{frame}

% ---- T2-F19 · EDIT · TIER: READ · from T2:653
\begin{frame}{Close the Loop: Reflect, Save, Shrink}
\textbf{The last five minutes of a session are worth more than the next fifty. Three habits.}

\vspace{0.2cm}

\begin{enumerate}
    \item \textbf{Reflect.} End with a two-minute review: what worked, what you corrected, what surprised you. Ask the agent itself to summarize the session's decisions --- it is cheap and surprisingly useful
    \medskip
    \item \textbf{Save.} Write that summary into the plan or progress file, and update \texttt{CLAUDE.md} while it is fresh. The next session then starts from the updated plan, not from anyone's memory
    \medskip
    \item \textbf{Shrink.} Prefer several small, focused sessions over one giant thread: a small context is cheaper, and the model's recall is sharper (mechanics in the next section). Break the big project into bounded sessions --- \emph{one intermediate product per session}
\end{enumerate}

\vspace{0.2cm}

\begin{takeawaybox}
\footnotesize
\textbf{A big project is a sequence of small sessions, each ending with an updated plan.} T5b turns this rhythm into a management system.
\end{takeawaybox}
\end{frame}


%=============================================================================
\section{Budget}
%===========================================================

\sectiondivider{Budget}{Tokens, cost, and the context window}

% ---- T2-F20 · KEEP · TIER: READ · from T2:681
\begin{frame}{Tokens and Cost on the Radar}
\textbf{Every agent turn is a model API call. Cost scales with input and output tokens.}

\vspace{0.2cm}

\begin{itemize}
    \item \textbf{Input tokens.} Your prompt, conversation history, \texttt{CLAUDE.md}, rules, retrieved chunks. Every step that adds context adds cost.
    \item \textbf{Output tokens.} Generated code, JSON tool calls, prose explanations. Reasoning-heavy modes also produce many internal tokens.
    \item \textbf{Long sessions accumulate context.} Watch the status line; check \texttt{/cost} before long tasks.
    \item \textbf{Pin for reproducibility.} Record the model version, prompt template, and seed when applicable; AEA disclosure (Oct 2024) extends to AI-assisted analysis.
\end{itemize}

\end{frame}

% ---- T2-F21 · KEEP · TIER: READ · from T2:695
\begin{frame}{The Context Window Is the Agent's Working Memory}
\textbf{Everything the model ``knows'' during a turn sits in one finite window (Lec07 T3b). Here is what fills it in an agent session:}

\vspace{0.15cm}

\begin{center}
\begin{tikzpicture}[
    seg/.style={draw=black!55, minimum height=0.72cm, font=\tiny, align=center, inner sep=1.5pt},
    wcap/.style={font=\tiny\itshape, text=black!65, align=center}
]
    \node[seg, fill=blue!12,  minimum width=1.9cm, text width=1.8cm] (s1) at (0,0) {system prompt + tool schemas};
    \node[seg, fill=blue!7,   minimum width=2.1cm, text width=2.0cm, right=0cm of s1] (s2) {\texttt{CLAUDE.md} + rules + auto memory};
    \node[seg, fill=blue!4,   minimum width=1.3cm, text width=1.2cm, right=0cm of s2] (s3) {skill index};
    \node[seg, fill=green!8,  minimum width=2.6cm, text width=2.5cm, right=0cm of s3] (s4) {conversation so far};
    \node[seg, fill=orange!12, minimum width=3.6cm, text width=3.5cm, right=0cm of s4] (s5) {tool outputs: file contents, command results, errors};
    \node[wcap] at ($(s1.south)!0.5!(s3.south)+(0,-0.32)$) {loaded at initialization (the frame before)};
    \node[wcap] at ($(s4.south)!0.5!(s5.south)+(0,-0.32)$) {grows every turn $\longrightarrow$ dominates long sessions};
\end{tikzpicture}
\end{center}

\vspace{0.2cm}

\begin{itemize}
    \item The window is a \textbf{finite budget}: model-dependent, but always exhaustible by a long session
    \item Recall is not uniform --- material buried mid-window is recalled worst (``lost in the middle,'' Lec07 T3b); we do not reteach it, we \emph{manage around} it
    \item Consequence: \textbf{long sessions degrade} --- the agent slowly forgets early decisions and corrections
\end{itemize}
\end{frame}

% ---- T2-F22 · EDIT · TIER: READ · from T2:724
% TODO(Phase B): Footnote the /compact <focus> syntax and the ~20-turn heuristic from archive v1 (Context Window Mechanics), tagged dated.
\begin{frame}{Compaction: What \texttt{/compact} Keeps and What It Drops}
\textbf{When the window fills, the harness summarizes: that is compaction. It is a lossy operation, and you should treat it as one.}

\vspace{0.1cm}

\begin{itemize}
    \item \texttt{/compact} = \textbf{summarize and clear}: older tool outputs are dropped first, the dialogue is condensed into a summary, and \texttt{CLAUDE.md} + rules are re-read from disk
    \item \textbf{What survives:} the decisions, the current task, key snippets. \textbf{What often does not:} the \emph{reasoning} behind decisions, rejected alternatives, small corrections you made in passing
    \item \textbf{Auto-compaction:} by default the harness compacts on its own as the window fills --- convenient, but the same information loss applies whether you triggered it or not
    \item \texttt{/clear} = \textbf{full reset}: conversation gone, files and memory remain; right choice when switching to an unrelated task
\end{itemize}

\vspace{0.1cm}

\begin{takeawaybox}
\footnotesize
\textbf{Rule: state belongs in files, not in the conversation.} If a decision matters tomorrow, write it to \texttt{CLAUDE.md} or a progress note \emph{before} the window fills --- then compaction costs you nothing.
\end{takeawaybox}

\vspace{0.05cm}
\footnotesize\textit{That is the mechanics. T5c returns to what this trade-off costs research integrity --- and why more context is not always better.}
\end{frame}

% ---- NEW · TIER: READ
% TODO(Phase B): Add the 9.3 economic twin (dividing the labour) in Phase B.
\begin{frame}{Bridge: From Using an Agent to Building One}
\textbf{You can install, launch, steer, budget, and read the screen. In the lab folder sit six agent files you have not yet written.}

\vspace{0.25cm}

\begin{itemize}
    \item \textbf{Next (Part 9.3):} build one from scratch (T3a), then learn the three ways to call it (T3b)
\end{itemize}
\end{frame}

%===========================================================
\end{document}
